\section{Methodology}
\label{sec:method}

Our goal is to characterize the correlation structure of LLM trustworthiness dimensions after removing the effect of two known confounds: model scale and RLHF alignment status. The key design principle is \textbf{confound removal before structure discovery}: raw Spearman correlations show mostly positive values because larger, better-trained models score higher on all dimensions --- masking the genuine trustworthiness geometry. Figure~\ref{fig:heatmaps} illustrates the difference between raw and partial correlation heatmaps.

\subsection{Data}

\textbf{Dataset.} We use TrustLLM \cite{sun2024trustllm} published evaluation scores: a 16-model $\times$ 6-dimension matrix with dimensions \{truthfulness (T), safety (S), fairness (F), adversarial robustness (R), privacy (P), machine\_ethics (E)\}. Scores are published as normalized values in $[0, 1]$.

\textbf{Models.} The 16-model set spans a range of scales ($\sim$7B to $\sim$175B parameters) and alignment approaches, including base models (LLaMA-2 7B/13B/70B, Falcon-7B, Mistral-7B) and RLHF-fine-tuned chat/instruct variants (LLaMA-2-7B-Chat, -13B-Chat, -70B-Chat, Vicuna-7B/13B, ChatGLM, Baichuan, WizardLM, Claude-2, GPT-3.5-turbo, GPT-4).

\textbf{VIF Check.} Variance inflation factors: VIF($\log_{10}$\_params) = 3.37, VIF(is\_RLHF) = 3.37 --- confirming acceptable collinearity ($< 5$) for partial correlation analysis with $df = 12$.

\textbf{Data provenance note.} All reported statistics are from the final validated experimental run documented in the Phase 4 synthesis report. Earlier exploratory runs during development used different implementation variants that may show variation in specific pair values; only the final validated outputs are reported here.

\subsection{Partial Spearman Correlation via OLS Residualization}

We compute partial Spearman correlations using OLS residualization:

\begin{algorithm}[h]
\caption{Partial Spearman via OLS Residualization}
\label{alg:partial_spearman}
\begin{algorithmic}[1]
\FOR{each dimension $d \in \{T, S, F, R, P, E\}$}
  \STATE $\text{rank}_d \leftarrow \texttt{scipy.stats.rankdata}(\text{scores}[:, d])$
  \STATE $\text{residual}_d \leftarrow \texttt{OLS}(\text{rank}_d \sim \log_{10}\text{params} + \text{is\_RLHF}).\text{resid}$
\ENDFOR
\FOR{each pair $(i, j)$}
  \STATE $\rho_\text{partial}(i,j) \leftarrow \texttt{scipy.stats.spearmanr}(\text{residual}_i, \text{residual}_j).\text{statistic}$
  \STATE $t \leftarrow \rho \sqrt{(n-2)/(1-\rho^2)}$
  \STATE $p \leftarrow 2 \cdot \texttt{scipy.stats.t.sf}(|t|, df=12)$
\ENDFOR
\end{algorithmic}
\end{algorithm}

The effective $df = 12$ accounts for 2 degrees of freedom consumed by OLS residualization ($n - k_\text{covariates} - 2 = 16 - 2 - 2 = 12$). A naive $t$-test on residuals would use $df=14$, but estimated (not true) residuals reduce effective $df$ to 12 (conservative).

\textbf{Bonferroni correction:} $\alpha_\text{corrected} = 0.05/15 = 0.0033$.

\subsection{Hierarchical Clustering}

We construct a distance matrix $d(i, j) = 1 - |\rho_\text{partial}(i, j)|$ and apply Ward, average, and complete linkage methods using \texttt{scipy.cluster.hierarchy.linkage} with precomputed distances. Cluster quality is assessed via silhouette score. \textbf{Implementation note:} sklearn Ward linkage does not support precomputed distance matrices (sklearn issue \#27655); we use scipy throughout.

\subsection{Minimum Spanning Tree and Evaluation Set}

The MST of the distance matrix identifies the minimum connected subgraph. Hub nodes (high MST degree) constitute the minimum evaluation set. Bootstrap stability is assessed via 1000 resamples of 14/16 models, reporting mean per-edge frequency \cite{tumminello2005tool} with threshold $\geq 0.90$.

\textbf{Metric choice.} We adopt the mean per-edge bootstrap frequency metric \cite{tumminello2005tool} rather than full-topology stability (Jaccard similarity of MST edge sets). The full-topology metric is sensitive to minor permutations in near-tie edge weights and yielded 0.606 in preliminary analysis --- before we recognized that edge-level frequency better characterizes stability for sparse trees with near-tie distances at small $n$. The mean per-edge metric is the standard in the financial correlation MST literature \cite{tumminello2005tool} and yields 0.917 in our final analysis.

\subsection{RLHF Mechanism: Within-Family Natural Experiment}

For LLaMA-2 (7B, 13B, 70B) base--chat pairs, we compute $\Delta_d = \text{score}_\text{Chat}(d) - \text{score}_\text{Base}(d)$ and test joint safety+ethics optimization via binomial sign test ($n = 3$). Figure~\ref{fig:scatter} illustrates OLS residualization for the safety--privacy pair.
